\documentclass[10pt,a4paper]{amsart}
\usepackage[margin=19mm]{geometry}
\usepackage{amsmath,amssymb,mathtools}
\usepackage[scr=rsfs,cal=euler]{mathalfa}
\usepackage{xfrac}
\usepackage[unicode,hidelinks,bookmarks=false]{hyperref}
\newcommand{\ie}{i.e.\ }
\newcommand{\cf}{cf.\ }
\newcommand{\resp}{respectively\ }
\newcommand{\Subsection}{\subsection}
\providecommand{\nobreakdash}{\nobreak\hbox{-}}
\setlength{\emergencystretch}{2em}
\multlinegap0pt
\usepackage{amssymb}
\usepackage{float}
\usepackage{xcolor}
\newtheorem{theorem}{Theorem}
\newtheorem{proposition}{Proposition}
\title{Restricted generalized Schur numbers}
\author{Collier Gaiser}
\date{Source: arXiv:2608.08789v1 (CC BY 4.0)}
\begin{document}
\maketitle

\noindent\emph{Source and licence.} Restricted generalized Schur numbers. Version arXiv:2608.08789v1 (CC BY 4.0), distributed by arXiv under the Creative Commons Attribution 4.0 International licence, http://creativecommons.org/licenses/by/4.0/, as recorded in the arXiv licence metadata for this exact version and shown on the abstract page https://arxiv.org/abs/2608.08789v1. Full licence notice: \texttt{licenses/arxiv-2608.08789-CC-BY-4.0.txt}.\par\smallskip

\noindent\textit{Editorial scope.} Counted unit: the article's proposition proposition (source0.tex lines 487--489). The article's complete original proof of it is delivered with it (source0.tex lines 490--505, one \texttt{proof} environment of 16 lines). No mathematics is written, repaired or re-derived: the statement and the proof are the article's own lines, in the article's own notation, and no retained line is substituted or re-flowed.  Retained uncounted as the source-local context the proof needs: lines 65--67; lines 82--92.  Nothing was omitted from any retained span, so the package carries no omission record.  No retained line contains a citation, so the wrapper carries no bibliography and no bibliography entry.  No retained line contains a figure environment, an image-inclusion command or drawing code, so no image is delivered and the package declares no asset dependency.  Editorial formatting only, in addition: an AMS \texttt{amsart} preamble that restates the article's own declarations and macros used by the retained lines, namely the environments \texttt{theorem}, \texttt{proposition} on separate counters, together with the standard packages those lines need.  The retained environments print in this document as Theorem 1, Proposition 1 in order of appearance, whereas the article prints the same items with the numbering of its own full document.  The complete native source of the article as distributed in the arXiv e-print for this exact version is bundled unchanged as \texttt{sources/207191/source0.tex}.
\begin{equation}\label{Equaiton:Main}
x_1+x_2+\cdots+x_k=x_{k+1}.
\end{equation}

\begin{theorem}\label{Theorem:Main}
For all $k,\ell\in\mathbb{N}$ with $2\leq\ell\leq k$, we have
    \[
    S_2(k;\ell)\geq k^2+\left[\frac{(\ell+1)(\ell-2)}{2}+2\right]k+\ell(\ell-2);
    \]
furthermore, fixing $\ell$, if $k$ is large enough, then we have
 \[
    S_2(k;\ell)= k^2+\left[\frac{(\ell+1)(\ell-2)}{2}+2\right]k+\ell(\ell-2).
 \]
In particular, we have $S_2(k;2)=k^2+2k$ for all $k\geq3$.
\end{theorem}

\begin{proposition}
We have $S_3(k;2)\geq k^3+3k^2+k-1$ for all $k\geq2$.
\end{proposition}

\begin{proof}
Let $k\geq2$ and let
\[
A_R=\{1,2,\ldots,k\}\cup\{k^2+k+1,k^2+k+2,\ldots,k^2+2k-1\}\cup\{k^3+2k^2+1,k^3+2k^2+2,\ldots,k^3+2k^2+k-1\},
\]
\[
A_B=\{k+1,k+2,\ldots,k^2+k\}\cup\{k^3+2k^2+k,k^3+2k^2+k+1,\ldots,k^3+3k^2+k-2\},
\]
and
\[
A_G=\{k^2+2k,k^2+2k+1,\ldots,k^3+2k^2\}.
\]
Let $\Delta:\{1,2,\ldots,k^3+3k^2+k-2\}\to\{R,B,G\}$ be a $3$-coloring such that $\Delta(a)=R$ for all $a\in A_R$, $\Delta(a)=B$ for all $a\in A_B$, and $\Delta(a)=G$ for all $a\in A_G$.

Similar to the proof of the lower bound in Theorem~\ref{Theorem:Main}, it is straightforward to check that none of $A_R$, $A_B$, or $A_G$ has a solution $\mathcal{S}$ to (\ref{Equaiton:Main}) such that $c(\mathcal{S})=3$. Hence $\Delta$ does not have a solution $\mathcal{S}$ to (\ref{Equaiton:Main}) such that $c(\mathcal{S})=3$. This completes the proof.
\end{proof}

\end{document}
